\section*{Related Work}\label{section:motivation}

The deployment of LLM-based clinical assistants in internet hospital settings presents a fundamental tension between the efficiency gains of automation and the personal nature of the doctor-patient relationship. In this section, we articulate the core motivations for doctor-specific personalization and situate our work within the broader research landscape.

\subsection*{The Personalized LLM Assistant in Clinical Communication}

Doctor-patient communication varies a lot from one doctor to another, for three main reasons. First, doctors differ in their language style and medical knowledge, both shaped by their specialty training and years of practice. A cardiologist talking about chest pain uses different words, tone, and structure than a dermatologist describing a skin condition, and a general AI model cannot pick up on these individual differences. Second, different hospital departments follow different communication rules: cardiology needs clear language about risks, while dermatology focuses on describing what things look like, and these department-specific norms are often invisible to general-purpose AI. Third, doctors are needed to be responsible for their reply to patients. If an AI suggestion does not sound like the doctor or goes against department rules, the doctor must spend a lot of time rewriting it, which defeats the point of using AI to improve efficiency. For these reasons, each doctor needs a personalized AI assistant that matches how they talk, follows their department's rules, and cuts down on manual rewriting~\cite{Johnson2021Precision}.

Existing approaches to clinical NLP largely ignore this dimension of personalization. Medical dialogue systems trained on aggregated corpora~\cite{zengMedDialogLargescaleMedical2020,moorFoundationModelsGeneralist2023} produce responses optimized for average performance across all clinicians, effectively erasing individual differences. While such models may achieve reasonable scores on automated metrics, their clinical utility is constrained by the ``last-mile'' problem of physician adoption~\cite{info:doi/10.2196/jmir.7854,rajpurkarAIHealthMedicine2022a}. Our work is motivated by the hypothesis that \textit{physician-specific model adaptation, when implemented efficiently, can close this adoption gap} by producing suggestions that physicians recognize as consistent with their own voice and practice norms.

\subsection*{Parameter-Efficient Fine-Tuning with LoRA}

Full fine-tuning of large language models for each individual physician is computationally prohibitive, particularly in hospital settings where dozens or hundreds of doctors may require personalized models. Low-Rank Adaptation (LoRA)~\cite{hu2021lora} addresses this challenge by freezing the pre-trained base model weights and injecting trainable low-rank decomposition matrices into the attention layers. This reduces the number of trainable parameters by several orders of magnitude while achieving competitive performance on downstream tasks. We adopt LoRA fine-tuning based on LLaMA-Factory~\cite{zheng2024llamafactory}, an efficient open-source framework that provides optimized kernels for training and inference.

Recent advances in parameter-efficient fine-tuning (PEFT)~\cite{peft} have introduced several adapter variants, including prefix tuning, prompt tuning, and IA$^3$. Among these, LoRA offers a favorable trade-off between adaptation quality and storage efficiency for our use case, as each doctor-specific adapter occupies only a few megabytes of storage and can be loaded and unloaded at inference time with minimal overhead.

\subsection*{LoRA Routing and Augmented Generation}

When multiple LoRA adapters coexist---each tailored to a different physician---the system must determine which adapter to activate for a given incoming query. LoRA composition methods such as LoRAHub~\cite{huangLoraHubEfficientCrossTask2024} dynamically select and compose LoRA modules for downstream tasks. In our setting, the routing problem is simplified by the fact that the target physician is known at query time (the patient is consulting a specific doctor), but dynamic adapter selection provides a principled framework for generalizing to scenarios where the optimal adapter must be inferred from the query context alone, such as triage or automated preliminary response generation.

\subsection*{Model Merging and Interference Resolution}

Deploying hundreds of individual LoRA adapters introduces operational complexity. An alternative is to merge multiple adapters into a single unified model, enabling a single inference endpoint to serve multiple physicians. However, naive parameter averaging often leads to destructive interference between adapters, where the merged model performs significantly worse than any individual adapter. TIES-MERGING~\cite{yadavTIESMergingResolvingInterference2023a} addresses this via three stages: (1) trimming low-magnitude parameter changes that may represent noise, (2) resolving sign conflicts through an elected sign mechanism, and (3) merging only the parameters that agree in sign direction. Our SCF method builds on this paradigm with enhanced sign-consistency mechanisms that further improve merging quality for large adapter pools. Task arithmetic~\cite{ilharco2023editing} provides an alternative formulation, but sign-resolved methods have been shown to outperform it when merging a larger number of task-specific models, making the approach particularly suitable for our multi-doctor scenario.



